%================================================================
\section{Replication}\label{app:replication}

Everything in the paper reruns with one command from the repository root, after a
one-time setup of the pinned environment described in \texttt{README.md}:
\begin{quote}
\texttt{.venv/bin/python scripts/replicate.py}
\end{quote}
The script runs the test suite, checks that the planted targets regenerate byte for
byte, runs the planted pipeline, the Foil~2 calibration check, the sensitivity menu
(amendment A1), the interference trade-off (amendment A2) and its descriptive
variance split (addendum A2.1), redraws the figures
from the saved results, copies them into \texttt{writeup/figures/} and
\texttt{thesis/figures/}, writes both copies of \texttt{numbers.tex}, builds the
earlier write-up, and compares every result file with the committed version. Every
random stream is seeded, so a rerun reproduces every digit; a mismatch means that
the code or the environment changed. The option \texttt{--quick} skips the Foil~2
check, both amendments and the addendum. This paper is then built from \texttt{thesis/} with
\texttt{latexmk}. The measured runtimes and package pins are in
\texttt{README.md}. The exogenous design lives in \texttt{config/world.toml} and
\texttt{config/design.toml}; the targets in \texttt{data/planted.json}; the results
in \texttt{results/}.

